% ===== 前置部分：标题页、摘要、译者说明、目录（Lead 独占）=====

\begin{titlepage}
\centering
\vspace*{2cm}
{\LARGE\bfseries 广义对称性讲义\par}
\vspace{0.5em}
{\large\itshape Lectures on Generalized Symmetries\par}
\vspace{2.2em}
{\large Lakshya Bhardwaj,\ Lea E. Bottini,\ Ludovic Fraser-Taliente,\\[2pt]
Liam Gladden,\ Dewi S. W. Gould,\ Arthur Platschorre,\ Hannah Tillim\par}
\vspace{1em}
{\small（据 Lakshya Bhardwaj 所授课程整理）\par}
\vspace{2em}
{\small arXiv:2307.07547v2 [hep-th]\quad 2023 年 9 月 8 日\par}
\vfill
\begin{minipage}{0.86\textwidth}
\small
\textbf{中译注本说明。}本文是上述讲义的完整中文译注本：正文逐段翻译，全部 745 条编号公式与原文
(2.1)--(6.54) 一一对应，图 1--28、表 1 与原文一致，参考文献保留英文原貌。
源码中以 \texttt{\% [Tutor Note: ...]} 形式写下的批注不进入排版，汇总见《译注汇总.md》。
\end{minipage}
\vspace{1.2em}
\end{titlepage}

\section*{摘要}
\addcontentsline{toc}{section}{摘要}

这是一套关于量子场论中\textbf{广义整体对称性}（generalized global symmetries）的讲义。讨论重点是可逆对称性，
并对不可逆对称性略作评论。主要论题包括高阶形式对称性的基本概念及其性质，特别是 't Hooft 反常、规范化（gauging）
与自发对称性破缺。我们还会介绍对称性拓扑场论（symmetry topological field theories，SymTFTs）这一有用概念。
此外，讲义给出高阶群对称性（描述高阶形式对称性之间混合的对称性）的导引。少数进阶论题包括全息（holography）
与弦论中的几何工程（geometric engineering）构造如何编码高阶形式对称性。全篇始终以规范理论为例来说明各项概念。

\vspace{1.2em}

\tableofcontents

\clearpage
